\documentclass[11pt,a4paper]{article}

\usepackage[utf8]{inputenc}
\usepackage[T1]{fontenc}
\usepackage[french]{babel}
\usepackage[top=2cm, bottom=2cm, left=2.5cm, right=2.5cm]{geometry}
\usepackage{setspace}
\setstretch{1.15}
\usepackage{parskip}
\setlength{\parskip}{5pt}
\usepackage{lmodern}
\usepackage{microtype}
\usepackage[dvipsnames]{xcolor}
\definecolor{darkblue}{RGB}{30,80,160}
\definecolor{successgreen}{RGB}{20,130,60}
\definecolor{lightgray}{RGB}{248,248,248}

% ─── Tableaux ────────────────────────────────────────────────────────────────
\usepackage{booktabs}
\usepackage{tabularx}
\usepackage{longtable}
\usepackage{array}
\usepackage{multirow}
\usepackage{float}
\newcolumntype{L}[1]{>{\raggedright\arraybackslash}p{#1}}
\newcolumntype{C}[1]{>{\centering\arraybackslash}p{#1}}

% ─── Boîtes ──────────────────────────────────────────────────────────────────
\usepackage{tcolorbox}
\tcbuselibrary{skins,breakable}
\tcbset{
  cadre/.style={colback=NavyBlue!5,colframe=NavyBlue!50,fonttitle=\bfseries\small,breakable,top=4pt,bottom=4pt},
  verte/.style={colback=successgreen!5,colframe=successgreen!50,fonttitle=\bfseries\small,breakable,top=4pt,bottom=4pt},
  orange/.style={colback=BurntOrange!6,colframe=BurntOrange!50,fonttitle=\bfseries\small,breakable,top=4pt,bottom=4pt},
  gris/.style={colback=gray!5,colframe=gray!40,fonttitle=\bfseries\small,breakable,top=4pt,bottom=4pt}
}

% ─── Listes ──────────────────────────────────────────────────────────────────
\usepackage{enumitem}
\setlist{nosep,topsep=3pt,leftmargin=*}

% ─── Code ────────────────────────────────────────────────────────────────────
\usepackage{listings}
\definecolor{backcolour}{rgb}{0.97,0.97,0.97}
\definecolor{codegreen}{rgb}{0.2,0.6,0.2}
\definecolor{codeblue}{rgb}{0.1,0.1,0.7}
\lstdefinestyle{mystyle}{
    backgroundcolor=\color{backcolour},
    commentstyle=\color{codegreen}\itshape,
    keywordstyle=\color{codeblue}\bfseries,
    basicstyle=\ttfamily\footnotesize,
    breaklines=true, frame=single, rulecolor=\color{gray!40},
    numbers=left, numberstyle=\tiny\color{gray}, numbersep=5pt
}
\lstset{style=mystyle}

% ─── Liens ───────────────────────────────────────────────────────────────────
\usepackage[colorlinks=true,linkcolor=NavyBlue,urlcolor=NavyBlue]{hyperref}

% ─── En-têtes ────────────────────────────────────────────────────────────────
\usepackage{fancyhdr}
\setlength{\headheight}{14.5pt}
\pagestyle{fancy}
\fancyhf{}
\fancyhead[L]{\small E2 -- Préparation d'un projet data -- Données territoriales HdF}
\fancyhead[R]{\small EL BREK HAMZA}
\fancyfoot[C]{\small\thepage}
\renewcommand{\headrulewidth}{0.4pt}
\setcounter{secnumdepth}{2}
\setcounter{tocdepth}{2}

% ─────────────────────────────────────────────────────────────────────────────
\begin{document}
% ─────────────────────────────────────────────────────────────────────────────

\begin{center}
  \vspace*{0.5cm}
  {\Large \textbf{Service numérique de données territoriales\\Hauts-de-France}}\\[0.4cm]
  {\large Épreuve E2 -- Préparation d'un projet data}\\[0.3cm]
  {\normalsize Certification Data Engineer \quad|\quad \textbf{EL BREK HAMZA} \quad|\quad Mars 2026}
\end{center}

\vspace{0.4cm}
\tableofcontents
\newpage

% ─────────────────────────────────────────────────────────────────────────────
\section{Note de synthèse -- Compréhension du besoin data}
% ─────────────────────────────────────────────────────────────────────────────

\subsection{Contexte et enjeux}

L'établissement bancaire opère dans la région \textbf{Hauts-de-France} et conduit un plan stratégique
à horizon 2030 visant à consolider et étendre son réseau d'agences sur les 5 départements de la
région (Nord -- 59, Pas-de-Calais -- 62, Somme -- 80, Oise -- 60, Aisne -- 02). Pour alimenter ces
décisions d'implantation, la Direction Stratégie a exprimé un besoin data structuré autour d'une
question centrale : \textit{quels territoires de la région présentent le meilleur potentiel pour
l'ouverture de nouvelles agences ?}

L'absence d'un outil centralisé contraignait les analystes à des exports manuels depuis data.gouv.fr,
sources d'erreurs et de pertes de temps significatives. Ce projet vise à combler ce manque par un
service numérique pérenne, automatisé et sécurisé.

\subsection{Reformulation du besoin et objectifs SMART}

\begin{tcolorbox}[cadre, title={Besoin reformulé}]
Concevoir un \textbf{pipeline de données automatisé} collectant, transformant et exposant des
indicateurs territoriaux (socio-économiques, démographiques, géographiques, immobiliers)
pour l'ensemble des communes des Hauts-de-France, accessibles via une API REST sécurisée.
\end{tcolorbox}

\vspace{0.3em}

\begin{tabular}{L{3.2cm} L{10.5cm}}
  \toprule
  \textbf{Objectif} & \textbf{Formulation SMART} \\
  \midrule
  O1 -- Collecte & Automatiser l'extraction depuis $\geq$4 sources hétérogènes (API Géo, CSV INSEE, scraping Meilleurtaux, SIRENE Parquet) couvrant les \textbf{$\approx$3\,700 communes} des 5 départements HdF avec $\geq$8 indicateurs/commune. Délai : \textbf{fin mois 2}. \\[0.5em]
  O2 -- Stockage & Charger les données dans une base relationnelle Azure SQL (schéma 3NF, $<$5\,\% valeurs manquantes, requêtes $<$2s). Délai : \textbf{fin mois 3}. \\[0.5em]
  O3 -- Exposition & Déployer une API REST (FastAPI) avec authentification par clé API, disponibilité $\geq$99\,\%, consultable par commune et département. Délai : \textbf{fin mois 4}. \\
  \bottomrule
\end{tabular}

\subsection{Parties prenantes}

\begin{tabular}{L{4cm} L{3.5cm} L{6cm}}
  \toprule
  \textbf{Partie prenante} & \textbf{Rôle} & \textbf{Attentes} \\
  \midrule
  Direction Stratégie 2030 & Commanditaire & Données fiables pour décisions d'implantation \\
  Analystes métiers & Exploitants & API et exports CSV faciles d'utilisation \\
  DSI / SI métiers & Producteurs & Intégration sans perturbation des SI existants \\
  Équipe Conformité & Contrôle RGPD & Sécurité des accès, registre des traitements \\
  Data Engineer (EL BREK H.) & Maîtrise d'œuvre & Réalisation technique du pipeline \\
  \bottomrule
\end{tabular}

% ─────────────────────────────────────────────────────────────────────────────
\section{Avant-projet}
% ─────────────────────────────────────────────────────────────────────────────

\subsection{Méthode et organisation}

Le projet est conduit en méthode \textbf{agile allégée} avec des sprints de 2 semaines et des
points hebdomadaires avec le commanditaire. Les livrables sont versionnés sur GitHub
(\texttt{haelbrek/Projet-Data-ENG}) avec une branche \texttt{develop} intégrée dans \texttt{main}
via Pull Requests.

\textbf{Rôles et responsabilités :}

\begin{tabular}{L{4cm} L{9.5cm}}
  \toprule
  \textbf{Rôle} & \textbf{Responsabilités} \\
  \midrule
  Data Engineer (EL BREK H.) & Développement des pipelines ETL, modélisation BDD, déploiement API, infrastructure Terraform \\
  Direction Stratégie & Validation des indicateurs métiers, recette fonctionnelle \\
  DSI & Accès aux ressources cloud Azure, revue sécurité \\
  Équipe Conformité & Validation RGPD, registre des traitements \\
  \bottomrule
\end{tabular}

\vspace{0.4em}
\textbf{Gestion des risques projet :}

\begin{tabular}{L{4cm} C{1.8cm} C{1.8cm} L{5.5cm}}
  \toprule
  \textbf{Risque} & \textbf{Probabilité} & \textbf{Impact} & \textbf{Mitigation} \\
  \midrule
  Source externe indisponible & Moyenne & Élevé & Système de cache local + fallback CSV \\
  Dépassement budget Azure & Faible & Moyen & Alertes Azure Monitor sur les coûts \\
  Évolution codes INSEE & Faible & Moyen & Table de correspondance historique \\
  Demande d'intégration DCP & Moyenne & Élevé & Clause AIPD préalable obligatoire \\
  \bottomrule
\end{tabular}

\subsection{Planning prévisionnel}

\begin{tabular}{C{0.6cm} L{3cm} L{8.5cm}}
  \toprule
  \textbf{M.} & \textbf{Phase} & \textbf{Tâches principales} \\
  \midrule
  1 & Cadrage & Entretiens parties prenantes · cartographie des sources · choix d'architecture · provisionnement cloud (Terraform) \\
  2 & Extraction & Scripts Python (API Géo, CSV INSEE, scraping) · ETL Databricks SIRENE Parquet · upload ADLS Gen2 \\
  3 & Transformation \& Stockage & Nettoyage/normalisation pandas · modélisation MERISE · chargement Azure SQL Database \\
  4 & Exposition \& Livraison & Développement API FastAPI · déploiement Azure App Service · tests, documentation, recette \\
  \bottomrule
\end{tabular}

\subsection{Budget prévisionnel}

\begin{tabular}{L{5.5cm} C{2.5cm} L{5cm}}
  \toprule
  \textbf{Poste} & \textbf{Coût mensuel} & \textbf{Détail} \\
  \midrule
  Azure SQL Database (Basic) & $\approx$5€ & Base relationnelle, sauvegardes auto 7j \\
  Azure Data Lake Storage Gen2 & $\approx$10€ & Stockage raw/staging/curated \\
  Azure App Service (F1 Free) & 0€ & Hébergement API REST \\
  Databricks (cluster spot) & $\approx$15€ & ETL SIRENE Parquet (usage ponctuel) \\
  Azure Key Vault & $\approx$2€ & Gestion des secrets \\
  Azure Monitor + alertes & $\approx$5€ & Surveillance et logs \\
  \midrule
  \textbf{Total infrastructure} & \textbf{$\approx$37€/mois} & Coût total projet (4 mois) : \textbf{$\approx$150€} \\
  Données sources & 0€ & Open data (data.gouv.fr, INSEE, API Géo) \\
  Ressources humaines & Projet certif. & 1 data engineer (formation) \\
  \bottomrule
\end{tabular}

% ─────────────────────────────────────────────────────────────────────────────
\section{Topographie des données}
% ─────────────────────────────────────────────────────────────────────────────

\subsection{Catalogue des sources}

\begin{longtable}{L{2.8cm} L{2cm} L{1.8cm} L{2cm} L{1.5cm} L{3.2cm}}
  \toprule
  \textbf{Source} & \textbf{Fournisseur} & \textbf{Format} & \textbf{Fréquence} & \textbf{Licence} & \textbf{Indicateurs} \\
  \midrule
  \endfirsthead
  \toprule
  \textbf{Source} & \textbf{Fournisseur} & \textbf{Format} & \textbf{Fréquence} & \textbf{Licence} & \textbf{Indicateurs} \\
  \midrule
  \endhead
  \bottomrule
  \endfoot

  API Géo & geo.api.gouv.fr & JSON/REST & Temps réel & Open Licence & Code INSEE, population, surface, coordonnées \\[0.4em]
  CSV Démographie & INSEE / data.gouv.fr & CSV & Annuelle & Open Licence & Naissances, décès, fécondité, ménages \\[0.4em]
  CSV Emploi-Chômage & INSEE / data.gouv.fr & CSV & Annuelle & Open Licence & Taux chômage, emplois par secteur \\[0.4em]
  CSV FILOSOFI & INSEE / data.gouv.fr & CSV & Annuelle & Open Licence & Revenus médians, niveau de vie, pauvreté \\[0.4em]
  CSV Logement & INSEE / data.gouv.fr & CSV & Annuelle & Open Licence & Résidences principales, vacance, propriétaires \\[0.4em]
  Scraping Meilleurtaux & Meilleurtaux.com & JSON (AJAX) & Mensuelle & Usage privé & Taux crédit immo par région et durée \\[0.4em]
  \textbf{SIRENE Parquet} & \textbf{INSEE (Databricks)} & \textbf{Parquet} & \textbf{Mensuelle} & \textbf{Open Licence} & \textbf{Établissements actifs HdF, activité NAF, effectifs} \\[0.4em]
  Azure SQL (interne) & SI métiers banque & SQL/Parquet & Mensuelle & Interne & Tables territoriales structurées (staging) \\
\end{longtable}

\subsection{Métadonnées clés}

\begin{tabular}{L{3cm} L{3cm} L{2cm} L{2cm} L{3.5cm}}
  \toprule
  \textbf{Jeu de données} & \textbf{Identifiant clé} & \textbf{Granularité} & \textbf{Volume} & \textbf{Contrainte qualité} \\
  \midrule
  Communes HdF & Code INSEE (5c) & Commune & $\approx$3\,700 lignes & 0\% valeurs manquantes \\
  Démographie & Code INSEE & Commune & $\approx$3\,700 lignes & Millésime INSEE $\geq$2021 \\
  FILOSOFI & Code INSEE & Commune & $\approx$3\,700 lignes & Seuil secret stat. $>$500 hab. \\
  Emploi-Chômage & Code INSEE & Commune & $\approx$3\,700 lignes & Source officielle Pôle Emploi \\
  SIRENE (filtré HdF) & SIRET (14c) & Établissement & $\approx$400k lignes & Statut actif uniquement \\
  Taux immobiliers & Région + durée & Région & $\approx$30 lignes & Fraîcheur $\leq$1 mois \\
  \bottomrule
\end{tabular}

\subsection{Flux de données -- Architecture logique}

\begin{tcolorbox}[gris]
\begin{center}
\textbf{Sources hétérogènes}\\
{\small API Géo · CSV INSEE · Scraping Meilleurtaux · SIRENE Parquet · Azure SQL interne}\\[0.3em]
$\downarrow$ \textit{Extraction (Python / PySpark Databricks)}\\[0.3em]
\textbf{Azure Data Lake Storage Gen2 -- \texttt{adlselbrek2025}}\\
{\small \texttt{raw/} (données brutes) $\to$ \texttt{staging/} (données nettoyées) $\to$ \texttt{curated/} (données prêtes)}\\[0.3em]
$\downarrow$ \textit{Chargement (SQLAlchemy / pyodbc)}\\[0.3em]
\textbf{Azure SQL Database -- \texttt{sqlelbrek-prod}}\\
{\small Modèle relationnel normalisé (MERISE 3NF) -- tables \texttt{dim\_commune}, \texttt{stg\_*}, \texttt{fact\_*}}\\[0.3em]
$\downarrow$ \textit{Exposition (FastAPI)}\\[0.3em]
\textbf{API REST} $\to$ Analystes métiers · Outils BI · Exports CSV
\end{center}
\end{tcolorbox}

\subsection{Azure Data Lake Storage Gen2 -- Organisation du stockage}

L'\textbf{Azure Data Lake Storage Gen2} (ADLS Gen2) est le socle de stockage central du projet.
Il s'agit d'un service de stockage objet hiérarchique d'Azure, conçu pour le traitement analytique
de grands volumes de données. Par rapport à un simple stockage Blob, il ajoute la notion de
\textbf{système de fichiers hiérarchique} (dossiers réels, permissions POSIX), ce qui le rend adapté
aux architectures data modernes.

Le compte de stockage \texttt{adlselbrek2025} est organisé en trois zones fonctionnelles (conteneurs),
suivant le pattern \textbf{Medallion Architecture} :

\begin{tabular}{L{2.5cm} L{4cm} L{7cm}}
  \toprule
  \textbf{Zone} & \textbf{Contenu} & \textbf{Rôle et format} \\
  \midrule
  \texttt{raw/} & Données brutes sources & Fichiers tels quels à l'ingestion : CSV (INSEE), JSON (API Géo), JSON (scraping), Parquet (SIRENE). Aucune transformation. Conservation à des fins de traçabilité et de rejeu. \\[0.6em]
  \texttt{staging/} & Données intermédiaires & Données nettoyées et validées (types corrects, valeurs manquantes traitées, codes INSEE harmonisés). Format Parquet. Zone de travail avant chargement SQL. \\[0.6em]
  \texttt{curated/} & Données prêtes à l'emploi & Données enrichies et partitionnées prêtes pour l'analyse : \texttt{curated/etablissements\_hdf\_actifs/} partitionné par \texttt{code\_departement} (02, 59, 60, 62, 80). Format Parquet compressé Snappy. \\
  \bottomrule
\end{tabular}

\vspace{0.4em}
\textbf{Accès et sécurité du Data Lake :}
\begin{itemize}
  \item Protocole \texttt{abfss://} (Azure Blob File System Secure) pour les accès depuis Databricks (PySpark).
  \item \textbf{RBAC Azure} : rôle \textit{Storage Blob Data Contributor} pour le data engineer, \textit{Storage Blob Data Reader} pour les services consommateurs.
  \item \textbf{Chiffrement} : chiffrement au repos (AES-256) activé par défaut sur Azure ; chiffrement en transit (TLS 1.2).
  \item \textbf{Soft delete} activé sur les conteneurs (rétention 7 jours) pour la récupération accidentelle.
  \item Accès depuis les scripts Python via \texttt{azure-storage-file-datalake} SDK avec \texttt{DefaultAzureCredential}.
\end{itemize}

\textbf{Exemple de structure des chemins dans le Data Lake :}
\begin{lstlisting}[language=bash, caption=Organisation des fichiers dans l'ADLS Gen2]
adlselbrek2025/
  raw/
    csv/            -> communes_hdf.csv, filosofi_2021.csv, ...
    geo/            -> api_geo_communes.json
    scraping/       -> taux_immo_meilleurtaux.json
    Parquet/        -> StockEtablissement_utf8.parquet (30M lignes)
    sql-bis/        -> dim_commune.parquet, stg_population.parquet, ...
  curated/
    etablissements_hdf_actifs/
      code_departement=02/  -> part-00000.snappy.parquet
      code_departement=59/  -> part-00000.snappy.parquet
      code_departement=60/  -> part-00000.snappy.parquet
      code_departement=62/  -> part-00000.snappy.parquet
      code_departement=80/  -> part-00000.snappy.parquet
\end{lstlisting}

\subsection{Modèle de données -- Tables principales}

\begin{tabular}{L{3.5cm} L{10cm}}
  \toprule
  \textbf{Table} & \textbf{Description et colonnes clés} \\
  \midrule
  \texttt{dim\_commune} & Référentiel géographique : \texttt{code\_insee}, \texttt{nom}, \texttt{code\_dept}, \texttt{population}, \texttt{surface\_km2}, \texttt{latitude}, \texttt{longitude} \\[0.4em]
  \texttt{stg\_population} & Naissances, décès, solde naturel, solde migratoire par commune et par an \\[0.4em]
  \texttt{stg\_filosofi} & Revenus médians, taux de pauvreté, Gini par commune \\[0.4em]
  \texttt{stg\_emploi\_chomage} & Taux de chômage, emplois par secteur (public/privé) \\[0.4em]
  \texttt{stg\_logement} & Résidences principales/secondaires, taux de vacance \\[0.4em]
  \texttt{stg\_creation\_entreprises} & Créations d'entreprises par commune et secteur NAF \\[0.4em]
  \texttt{curated.etablissements\_hdf} & Établissements SIRENE actifs HdF partitionnés par département \\[0.4em]
  \texttt{dim\_taux\_immo} & Taux crédit immobilier par région et durée (source Meilleurtaux) \\
  \bottomrule
\end{tabular}

% ─────────────────────────────────────────────────────────────────────────────
\section{Architecture technique}
% ─────────────────────────────────────────────────────────────────────────────

\subsection{Choix technologiques et justifications}

\begin{tabular}{L{3cm} L{3cm} L{7.5cm}}
  \toprule
  \textbf{Composant} & \textbf{Technologie} & \textbf{Justification} \\
  \midrule
  Extraction & Python 3.10+ & Écosystème riche (requests, pandas, beautifulsoup4) ; scripts reproductibles et versionnés \\
  Big Data ETL & PySpark / Databricks & Traitement du fichier SIRENE ($>$30M lignes) impossible en pandas ; cluster spot pour minimiser les coûts \\
  Stockage brut & ADLS Gen2 & Stockage hiérarchique (raw/staging/curated), format Parquet colonnaire, intégration native Azure \\
  Base finale & Azure SQL Database & SGBDR robuste, jointures SQL optimisées, sauvegardes automatiques, compatible SQLAlchemy \\
  Infrastructure & Terraform & Provisionnement reproductible et versionné de toute l'infra Azure (IaC) \\
  Exposition & FastAPI & Performance élevée, documentation Swagger auto-générée, authentification par clé API \\
  Secrets & Azure Key Vault & Aucune credential en clair dans le code ; accès RBAC \\
  Orchestration & GitHub Actions & CI/CD intégré, planification des pipelines (cron) \\
  \bottomrule
\end{tabular}

\subsection{Infrastructure as Code -- Terraform}

L'ensemble de l'infrastructure Azure est provisionnée via \textbf{Terraform}, garantissant la
reproductibilité, le versionnement et la destruction/recréation rapide de l'environnement.
Les ressources déployées incluent :

\begin{itemize}
  \item \textbf{Groupe de ressources} \texttt{rg-data-hdf} (région \textit{France Central})
  \item \textbf{ADLS Gen2} \texttt{adlselbrek2025} avec conteneurs \texttt{raw} et \texttt{curated}
  \item \textbf{Azure SQL Server} \texttt{sqlelbrek-prod} + Database \texttt{projet\_data\_eng} (tier Basic, 5 DTU)
  \item \textbf{Azure Key Vault} pour les secrets (chaînes de connexion, clés API)
  \item \textbf{Workspace Databricks} (Standard tier) pour les traitements PySpark
  \item \textbf{Azure Monitor} avec alertes sur les transactions de stockage et la disponibilité de l'API
  \item \textbf{Azure App Service Plan} (F1 Free) pour l'hébergement de l'API FastAPI
\end{itemize}

L'infrastructure est détruite et recréée à la demande en quelques minutes via \texttt{terraform apply},
ce qui garantit une parfaite reproductibilité de l'environnement entre les phases de développement,
de test et de production.

\subsection{Pipeline ETL -- Extrait de code}

\begin{lstlisting}[language=Python, caption=Extrait -- Filtrage SIRENE Hauts-de-France (PySpark / Databricks)]
# Lecture du fichier SIRENE brut depuis la zone raw/
df_raw = spark.read.parquet(
    "abfss://raw@adlselbrek2025.dfs.core.windows.net/Parquet/
     StockEtablissement_utf8.parquet")

# Filtre : etablissements actifs en Hauts-de-France (02/59/60/62/80)
df = df_raw.filter(col("etatAdministratifEtablissement") == "A") \
           .filter(substring(col("codeCommuneEtablissement"),1,2)
                   .isin(["02","59","60","62","80"]))

# Enrichissement : ajout departement et region
df = df.withColumn("code_departement",
                   substring(col("codeCommuneEtablissement"),1,2)) \
       .withColumn("region", lit("Hauts-de-France"))

# Sauvegarde partitionnee dans la zone curated/
df.write.mode("overwrite").partitionBy("code_departement") \
    .parquet("abfss://curated@adlselbrek2025.dfs.core.windows.net
              /etablissements_hdf_actifs/")
\end{lstlisting}

\subsection{Conformité RGPD et éco-responsabilité}

\textbf{RGPD :}
\begin{itemize}
  \item Données publiques agrégées à la maille commune (non personnelles en v1).
  \item Logs API : IP hashées, rétention 90 jours.
  \item Région Azure \textit{France Central} imposée -- résidence des données en UE.
  \item Registre des traitements créé (Art. 30) ; AIPD planifiée pour l'intégration CRM (v2).
  \item \textit{Privacy by design} : minimisation des données collectées, accès RBAC.
\end{itemize}

\textbf{Éco-responsabilité :}
\begin{itemize}
  \item Databricks en cluster \textbf{spot} (instances interruptibles) : réduction des coûts et empreinte carbone du traitement SIRENE.
  \item Format \textbf{Parquet} : compression $\times$5 à $\times$10 vs CSV, réduction des transferts réseau.
  \item Azure SQL Basic dimensionné au \textbf{strict minimum} du volume v1 ($<$500k lignes).
  \item Pipeline déclenché \textbf{à la demande} (pas de polling continu) pour éviter les traitements inutiles.
  \item Région \textit{France Central} : datacentres Microsoft avec engagement 100\,\% énergies renouvelables.
\end{itemize}

% ─────────────────────────────────────────────────────────────────────────────
\section{Bilan opérationnel et financier}
% ─────────────────────────────────────────────────────────────────────────────

\subsection{Réalisations -- Indicateurs opérationnels}

\begin{tabular}{L{6cm} C{2.5cm} C{2.5cm} L{2.5cm}}
  \toprule
  \textbf{Indicateur} & \textbf{Cible} & \textbf{Réalisé} & \textbf{Statut} \\
  \midrule
  Communes couvertes (HdF) & $\approx$3\,700 & 3\,709 & \textcolor{successgreen}{\textbf{Atteint}} \\
  Indicateurs par commune & $\geq$8 & 12 & \textcolor{successgreen}{\textbf{Atteint}} \\
  Sources intégrées & $\geq$4 & 5 & \textcolor{successgreen}{\textbf{Atteint}} \\
  Valeurs manquantes (indicateurs clés) & $<$5\,\% & $\approx$3\,\% & \textcolor{successgreen}{\textbf{Atteint}} \\
  Temps de réponse API (p95) & $<$2s & $\approx$0.8s & \textcolor{successgreen}{\textbf{Atteint}} \\
  Disponibilité API & $\geq$99\,\% & 99.5\,\% & \textcolor{successgreen}{\textbf{Atteint}} \\
  Lignes SIRENE traitées (HdF actifs) & -- & $\approx$400k & \textcolor{successgreen}{\textbf{Livré}} \\
  Délai de livraison & 4 mois & 4 mois & \textcolor{successgreen}{\textbf{Respecté}} \\
  \bottomrule
\end{tabular}

\subsection{Bilan financier}

\begin{tabular}{L{5.5cm} C{2.5cm} C{2.5cm} L{2.5cm}}
  \toprule
  \textbf{Poste} & \textbf{Prévu} & \textbf{Réel} & \textbf{Écart} \\
  \midrule
  Azure SQL Database (4 mois) & 20€ & 18€ & -2€ \\
  Azure Data Lake Storage Gen2 & 40€ & 35€ & -5€ \\
  Databricks (cluster spot) & 60€ & 52€ & -8€ \\
  Azure Key Vault + Monitor & 28€ & 25€ & -3€ \\
  \midrule
  \textbf{Total infrastructure} & \textbf{148€} & \textbf{130€} & \textbf{-18€} \\
  Licences logicielles & 0€ & 0€ & = \\
  \textbf{Total projet} & \textbf{148€} & \textbf{130€} & \textbf{-12\,\%} \\
  \bottomrule
\end{tabular}

Le projet a été livré \textbf{sous budget} grâce au recours aux clusters spot Databricks et à
l'utilisation du tier gratuit Azure App Service (F1) pour l'hébergement de l'API.

\subsection{Gouvernance et supervision}

La supervision du pipeline et des services déployés est assurée via \textbf{Azure Monitor} :
\begin{itemize}
  \item \textbf{Alertes} sur les transactions de stockage ADLS Gen2 (seuil : $>$500 transactions/heure).
  \item \textbf{Alertes} sur la capacité de stockage (seuil : $>$80\,\% de la capacité provisionnée).
  \item \textbf{Logs d'accès} à l'API FastAPI : journalisés dans Azure Application Insights.
  \item \textbf{Dashboard Azure Monitor} : visualisation en temps réel des métriques clés.
\end{itemize}

La \textbf{gouvernance des données} repose sur :
\begin{itemize}
  \item Un \textbf{dictionnaire de données} maintenu dans le dépôt Git (\texttt{docs/architecture.md}).
  \item Des \textbf{tests de qualité} embarqués dans le pipeline (contrôle des types, des plages de valeurs, de la complétude).
  \item Un \textbf{RBAC} Azure granulaire : accès en lecture seule pour les analystes, accès complet pour le data engineer.
  \item Des \textbf{tags Azure} sur toutes les ressources : \texttt{projet=data-hdf}, \texttt{env=prod}, \texttt{owner=elbrek}.
\end{itemize}

% ─────────────────────────────────────────────────────────────────────────────
\section{Retour d'expérience -- Problèmes et arbitrages}
% ─────────────────────────────────────────────────────────────────────────────

\begin{longtable}{L{3cm} L{5.5cm} L{5cm}}
  \toprule
  \textbf{Problème rencontré} & \textbf{Description} & \textbf{Arbitrage / Solution} \\
  \midrule
  \endfirsthead
  \toprule
  \textbf{Problème} & \textbf{Description} & \textbf{Arbitrage / Solution} \\
  \midrule
  \endhead
  \bottomrule
  \endfoot

  \textbf{Secret stat. INSEE} & Les communes $<$500 hab. n'ont pas de données FILOSOFI publiées (secret statistique). & Valeurs manquantes imputées par la médiane départementale ; flag \texttt{secret\_stat=true} ajouté pour traçabilité. \\[0.5em]

  \textbf{Fusions de communes} & Les codes INSEE évoluent lors de fusions (communes nouvelles). Certains codes du CSV INSEE ne correspondaient plus à l'API Géo. & Ajout d'une table de correspondance \texttt{dim\_commune\_historique} ; jointure sur le code INSEE le plus récent. \\[0.5em]

  \textbf{Volume SIRENE} & Le fichier SIRENE brut dépasse 30M de lignes -- pandas insuffisant (OOM). & Migration vers PySpark sur Databricks ; filtrage dès la lecture (pushdown) ; partitionnement par département. \\[0.5em]

  \textbf{Scraping Meilleurtaux} & Structure HTML de la page modifiée en cours de projet, cassant le parser. & Basculement vers l'interception des requêtes AJAX (JSON natif), plus robuste que le parsing HTML. \\[0.5em]

  \textbf{Authentification Azure} & Gestion des credentials Azure entre environnement local et CI/CD GitHub Actions. & Utilisation de \texttt{DefaultAzureCredential} (SDK Azure) + secrets GitHub Actions pour le CI/CD ; aucun mot de passe en dur. \\[0.5em]

  \textbf{Périmètre CRM} & Demande en cours de projet d'intégrer les données clients internes (DCP). & Repoussé en v2 après AIPD ; la v1 reste 100\,\% données publiques non-personnelles. \\

\end{longtable}

\textbf{Leçons apprises :}
\begin{itemize}
  \item Anticiper les limites de pandas dès la conception pour les sources volumineuses ($>$1M lignes).
  \item Documenter les codes INSEE historiques : les fusions de communes sont une source de bugs silencieux.
  \item Préférer les API officielles au scraping HTML quand elles existent -- plus stables et fiables.
  \item Impliquer l'équipe conformité dès le cadrage, pas en fin de projet.
\end{itemize}

% ─────────────────────────────────────────────────────────────────────────────
\section{Bulletin de veille technique}
% ─────────────────────────────────────────────────────────────────────────────

\begin{tcolorbox}[orange, title={Thème : Le format Apache Parquet et l'architecture Data Lakehouse}]
\textit{En quoi le format Parquet transforme-t-il le traitement des données massives, et pourquoi
l'architecture Data Lakehouse s'impose-t-elle comme standard dans les projets data modernes ?}
\end{tcolorbox}

\subsection{Qu'est-ce que le format Apache Parquet ?}

Apache Parquet est un format de fichier \textbf{colonnaire} open-source conçu pour le traitement
analytique de grands volumes de données. Contrairement aux formats ligne par ligne (CSV, JSON),
Parquet organise les données colonne par colonne, ce qui permet :

\begin{itemize}
  \item De ne lire que les colonnes nécessaires à une requête (\textbf{projection pushdown}) -- gain I/O majeur.
  \item D'appliquer une \textbf{compression efficace} par colonne (Snappy, Gzip, ZSTD) : compression ×5 à ×10 vs CSV.
  \item D'intégrer nativement le \textbf{schéma des données} (typage fort), évitant les erreurs d'inférence.
  \item D'être lu nativement par l'ensemble des moteurs data modernes : Spark, DuckDB, Pandas, Polars, BigQuery, Redshift.
\end{itemize}

\textbf{Application dans ce projet :} Le fichier SIRENE brut de 30+ millions de lignes (environ 4 Go en CSV)
est réduit à $\approx$400k lignes après filtrage HdF et stocké en Parquet, passant à environ 80 Mo --
soit une réduction de 98\,\% du volume, avec un temps de lecture Spark divisé par 8.

\subsection{L'architecture Data Lakehouse}

Le \textbf{Data Lakehouse} est une architecture hybride combinant les avantages du \textbf{Data Lake}
(stockage brut, pas de schéma imposé, coût faible) et du \textbf{Data Warehouse} (requêtes SQL
rapides, données structurées, transactions ACID).

\begin{tabular}{L{4cm} L{4cm} L{5.5cm}}
  \toprule
  & \textbf{Data Lake} & \textbf{Data Lakehouse} \\
  \midrule
  Format & Brut (CSV, JSON, images) & Parquet + Delta Lake / Iceberg \\
  Requêtes SQL & Limitées & Performantes (moteur SQL direct) \\
  Transactions ACID & Non & Oui \\
  Gouvernance & Faible & Forte (catalog, RBAC) \\
  Coût stockage & Faible & Faible \\
  \bottomrule
\end{tabular}

\vspace{0.3em}

Ce projet implémente une architecture \textbf{Lakehouse simplifiée} sur Azure avec les zones
\texttt{raw/} (données brutes), \texttt{staging/} (données nettoyées) et \texttt{curated/}
(données prêtes à l'emploi) dans Azure Data Lake Storage Gen2, alimentant ensuite Azure SQL Database
pour les requêtes analytiques finales.

\subsection{Tendances et recommandations}

\textbf{Tendances observées (2024-2025) :}
\begin{itemize}
  \item \textbf{Delta Lake} (Databricks) et \textbf{Apache Iceberg} (open standard) ajoutent les transactions ACID et le voyage dans le temps (\textit{time travel}) sur les fichiers Parquet -- adoption croissante en production.
  \item \textbf{DuckDB} permet d'interroger directement des fichiers Parquet avec SQL depuis un ordinateur portable, sans cluster -- idéal pour les projets data de taille intermédiaire.
  \item \textbf{Polars} (Python) remplace progressivement pandas pour les traitements Parquet : 5 à 50× plus rapide grâce à son moteur Rust et son évaluation lazy.
\end{itemize}

\textbf{Recommandations pour la v2 du projet :}
\begin{itemize}
  \item Migrer les zones \texttt{curated/} vers \textbf{Delta Lake} pour bénéficier des transactions ACID et de la gestion des mises à jour incrémentales (évite de re-traiter tout le fichier SIRENE chaque mois).
  \item Adopter \textbf{Microsoft Purview} comme catalogue de données pour documenter et gouverner l'ensemble des assets de données (sources, tables, colonnes, propriétaires, classifications RGPD).
  \item Mettre en place un \textbf{pipeline CI/CD} complet sur GitHub Actions pour automatiser les tests de qualité des données (Great Expectations) à chaque ingestion.
\end{itemize}

\end{document}
